\documentclass[a4paper, 12pt]{article}


\usepackage[utf8x]{inputenc}
\usepackage{fancyhdr}
\usepackage{lastpage}
% For figures and stuff
\usepackage{graphicx, wrapfig, subcaption, setspace, booktabs}
\usepackage[T1]{fontenc}
% Change for different font sizes and families
\usepackage[font=large, labelfont=bf]{caption}
\usepackage{fourier}
\usepackage[]{microtype}

\usepackage{graphicx}
\usepackage{float}
\usepackage{listings}
\usepackage{xcolor}
\usepackage{graphicx}

\usepackage[sorting=none]{biblatex} %Imports biblatex package
\addbibresource{sample.bib} %Import the bibliography file


\definecolor{codegreen}{rgb}{0,0.6,0}
\definecolor{codegray}{rgb}{0.5,0.5,0.5}
\definecolor{codepurple}{rgb}{0.58,0,0.82}
\definecolor{backcolour}{rgb}{0.95,0.95,0.92}

\lstdefinestyle{mystyle}{
    backgroundcolor=\color{backcolour},   
    commentstyle=\color{codegreen},
    keywordstyle=\color{magenta},
    numberstyle=\tiny\color{codegray},
    stringstyle=\color{codepurple},
    basicstyle=\ttfamily\normalsize,
    breakatwhitespace=false,         
    breaklines=true,    captionpos=b,                    
    keepspaces=true, numbers=left,                    
    numbersep=5pt, showspaces=false,                
    showstringspaces=false, showtabs=false, tabsize=2}
\lstset{style=mystyle}




\newcommand{\HRule}[1]{\rule{\linewidth}{#1}}
\onehalfspacing
\setcounter{tocdepth}{5}
\setcounter{secnumdepth}{5}

\usepackage[a4paper,top=2.5cm,bottom=2cm,left=2.5cm,right=2.5cm,marginparwidth=1.8cm]{geometry}
\pagestyle{fancy}


\setlength\headheight{5pt}
\fancyhead[L]{} 
\fancyhead[R]{}

\linespread{1.6}


%-----------------------------------------BEGIN DOC----------------------------------------
\begin{document}

\title{{\Huge \textbf{Literature Review}{\large\linebreak\\}}{\Large \textit{Criticism of Stakeholder capitalism model and respective ESG factor implementation. Brown \& Cohen (2023)\cite{BrownCohen}}\linebreak\linebreak}}

\author{Hanshuo Ma
\\ \\ \\
\textbf{University of Toronto Missisauga} \\
Department of Economics \\ \\
Professor Varouj A. Aivazian \\
ECO440: Corporate Finance \\ \\}
\date{\today}

\maketitle

\newpage
%-----------------------------------------CONTENT-------------------------------------
\tableofcontents\label{c}
\newpage

%--------------------------------------Introduction---------------------------------
\section{Summary of Stakeholder Capitalism and ESG factors} \label{Summary} 
%This paper is a literature review/critic review of the paper , which Brown and Cohen introduce Stakeholder Capitalism model and how does it works in  both theory and practice. The model incorporate ESG factors in corporate evaluation to perform better in long run compare to Shareholder primacy model in terms of managers able to maximize higher value of corporation despite making lower profits. 
%-----------------------------------------Part 1------------------------------------
\vspace{-10pt}
\subsection{The idea behind: Stakeholder vs Shareholder}
\vspace{-10pt}
  \hspace{12pt} Ever since Milton Friedman championed the philosophy of setting social consciences and corporate responsibility completely outside of the premise of corporate decision-making \cite{Milton}. The shareholder primacy model remains a highly debated and controversial model. Economists take this model for granted because its inherited assumption of profit-maximizing and efficient allocation of resources make modeling less complicated. In practice, the universal application of this model is also reflected in the market; Not only shareholder primacy model dominate U.S. discussions of corporate purpose, but conversations abroad as well. (Lynn Stout  2012)\cite{Lynn} However, in practice, the shareholder primacy model has its fair share of trouble and flaws. First of all, the inevitable agency cost between shareholders and debt-holders in a levered firm, or between shareholders and managers in an unleavened firm. Both scenarios will result in the company losing economic growth as the cost of the conflicts. Jensen and Meclking argued that this agency cost is inevitable because of the nature of the shareholder primacy model - separation of ownership and control.\cite{Jensen} Secondly, short-term profit-maximizing usually builds upon the opportunity cost of long-term growth, like massive layoffs and budget cuts across departments. Last of all and most important, a company that prioritizes shareholder primacy doesn't hold itself accountable for damaging public goods, such as air pollution, and river and ocean sewage pollution. 
 
 All the negative externalities the corporation created should be held accountable by the regulators, moreover, during the process of management decision-making, employees, and shareholders should be aware of the following consequences. Thus the premise of the Stakeholder model is to incorporate the demands of various change agents, including, employees, managers, consumers, suppliers, and regulators, as well as its shareholders - so the corporation could be working towards long-term value maximization for all of its stakeholders. 

 Stakeholders in general can be classified into three categories: (1) internal direct interaction with the company, (2) external direct interaction with the company, and (3) external indirect interaction helps define business operating environments. \cite{McKinsey} Nonetheless, each corporation has a unique composition of stakeholders and the corporation needs to conduct its research and analysis to identify and evaluate all stakeholder shares according to the corporation's strength and structure.  

\vspace{-10pt}
\subsection{Stakeholder in practice: ESG}
\vspace{-10pt}
  \hspace{12pt} In practice, the Stakeholder model encourages managers to balance many aspects of the company without resorting to trade-offs, there are three main pillars of balance, Environmental, Social, and (Corporate) Governance. For instance, balance sustainability and consumerism, employee satisfaction and high productivity, corporation economic growth and employee income disparities, and long-term sustainability against short-term profitability. However, since there isn't a standard practicable model that involves all stakeholders' interests in the decision-making process within corporate governance, one important underlying assumption is, to allow the outside investors and shareholders all have unanimous non-pecuniary preferences (i.e., non-monetary motivations) to incentive managers to consider ESG factors. On top of that, the model expands using the findings of Pastor, Stambaugh, and Taylor (2021, henceforth PST)\cite{Pastor}, whereas investors only consider how green (sustainable) the corporation is. 

  PST model finds that investors' preferences over the sustainability factor move asset pricing, the green stock has negative alphas - lower expected returns because the investors would like to hold them long-term to hedge against climate risk. And if at any time the investors and consumers are concerned about future sustainability or climate risk the green stock will outperform the non-green one in that period. PST model also finds that investing in the sustainability factor leads to positive social impact, whereby the corporation will have a higher market valuation due to investors willing to pay for a green premium and thus a lower cost of capital. 

  Here, not only the single sustainability factor in the PST model is expanded to $n$ numbers of arbitrarily equally important ESG factors, but the assumption of only the investor-related preferences driven upmarket evaluation is relaxed as well since now the managers are actively trying to invest more in projects that increase ESG factors rating within the corporation. The same expected future earnings with a higher evaluation of the company implies a lower cost of capital. Companies have more range and capital to choose for future projects since lower cost of capital makes some otherwise financially unviable projects viable. 

\vspace{-10pt}
\subsection{The future, and ESG investing}
\vspace{-10pt}
  \hspace{12pt} The strong non-pecuniary preferences of investors and shareholders advocate the company to disclose any ESG-related factors improved or invested within the corporation. Not only investors can hedge risk within their portfolio, but their investment also provides a lower cost of capital for a wider range of future projects the company couldn't finance otherwise. In contrast, investors with pecuniary preferences will invest in corporations with less commitment toward their ESG rating factors, and the stakeholders, of these corporations will be expected to generate a higher financial return. 

  Lower cost of capital, Lower cost of equity, and potential access to lower cost of debt - Green Bond investors favor corporations with higher commitment towards ESG rating factors, and having access to lower interest rates of debt could contribute towards potential change within the corporate capital structure. 


%-----------------------------------------Analysis-------------------------------------
\section{Criticism of Stakeholder Capitalism and ESG factors} \label{Criticism}
\vspace{-10pt}
\subsection{Minimal deviation from the shareholder model}\label{nonpecuniary preferences}
\vspace{-10pt}

  \hspace{12pt} The Stakeholder model in theory should have a solution to inherited agency cost in the Shareholder primacy model, however, in practice, the separation of ownership and control persists. Especially, now the model is the same shareholder model wrapped in a layer of ESG factors with investor and consumer preferences to incentivize managers to push the corporation in a relative direction. 
  
  The model's dependence on the non-pecuniary preferences of the shareholders and excluding the non-managerial stakeholder members from the corporation decision-making stage is inconsistent with the central idea behind Stakeholder capitalism. First of all, unanimous non-pecuniary preferences among the shareholders don't necessarily mean maximizing the long-term stability of the workforce and long-term consumers' interests. Investors' non-pecuniary preferences may be an extension of risk-averse preference hedging against climate risk rather than emphasizing the (S) Social and (G) Corporate governance factors in the corporation. In this scenario even as the long-term environmental impact of the corporation may be positive, the competence of corporate governance and its respective social impact on the local community remain debatable. Secondly, without unanimous non-pecuniary preferences, the stakeholder model will face the same agency problem as the shareholder primacy model does, and a trade-off between short-term profit and long-term sustainability will spark more uneconomic growth in the corporation. Even if all the investors have non-pecuniary preferences, if the shareholders possess different priorities relating to ESG factors, the manager may still forced to make decisions that only maximize the short-term profit of the corporation. Lastly, the PST model builds upon the risk dispersion between investors, if all the investors in the market hold sustainability-related products, everyone is risk-averse, and such low dispersion within the preference will dissipate the ESG industry.  


\subsection{Environment outweighs Social and Corporate Governance}
\vspace{-10pt}

  \hspace{12pt} Ultimately it could be the cost of investing in the Environmental factor is cheaper than the respective Social and Corporate Governance factors, we still observed that despite investors' tolerance towards long-term sustainability and lower expected return, the corporation still relies on massive layoffs and cancellation of future projects when facing financial turbulence or less than expected profitability in the short term. Recently, the massive layoffs in Meta, Google, Microsoft, and Hasbro all testify that the workforce is still burdened by short-term poor management in corporate government, and they are all still considered the leading companies in their industry in terms of ESG score and ESG commitment. 
  
  %Two reasons first could channel this decision, the utility of higher short-term profitability outweighs the higher long-term sustainability and thus loss of current corporate valuation due to shareholders' preferences; second, the ESG factors of environmental outweigh both factors of Social and Corporate Governance combined, which render the model closely to shareholder primacy model. Corporate governance should be held more accountable and incentivized to focus on long-term sustainability rather than sacrificing the workforce and compromising future economic growth, whenever the manager incorrectly predicts the short-term profitability and long-term growth opportunity in the industry.  

\subsection{ESG, Geopolitics, and future challenges}\label{ESG, Geopolitics, and future challenges}

\vspace{-5pt}
\subsubsection{Human rights in disputed territories}\label{Human rights in disputed territories}
\vspace{-10pt}
\hspace{12pt} International mega corporations are no strangers to human rights violations, operating in foreign regions with local management and operating within local law with a lack of international oversight and regulations can often lead to exploits of the local workforce. This is why have an international soft-law and human rights framework that was designed to compensate for corporate governance shortfall and ensure that corporations are held liable for their actions that impact the rights of others. Nonetheless, there is still empirical evidence suggesting that the corporation fails to meet its human rights responsibilities, especially operating in disputed territories, under the existing international framework. \cite{Ascencio}

The same international framework that the corporations are bound to, is forcing the corporations to make decisions that they aren't capable of making, nor they are sufficiently informed to make such decisions. The cost of conflict within corporate governance and the opportunity cost of losing rankings of ESG rating factors because of geopolitical challenges will often let the corporations decide to withdraw and not invest in the disputed territories. Discouraged corporations and lack of international support will often deny the potential economic growth of the disputed territory, furthermore, poverty creates conditions that deepen human rights violations and labor exploits which undermines the very purpose of the existing international law and human rights framework. 

\vspace{-5pt}
\subsubsection{Defense Industry and ESG contradiction}\label{Defense Industry and ESG contradiction}
\vspace{-10pt}
\hspace{12pt} The Defensive industry has been long viewed as a "sin" industry, thus the industry always has difficulty accessing financing. However, ever since the Russia-Ukraine war in 2022, defensive contractors in Europe have been met with a sudden change of attitude from the banks and the regulators - they are now considered the defenders of democracy and thus the defensive industry is classified as an ESG investment. (Causevic \& Causevic 2022)\cite{Causevic} However, it's easy to see the underlying contradiction within such classification, weapons, military vehicles, and missiles, are extremely bad for the environment, not to mention the amount of urban and environmental destruction that comes along with military actions. Most importantly, and defensive industry has been exploiting such a problem for decades- weapons don't end up where you send them. The social unrest and potentially increasing growth in terrorism are fueled by such an increased investment in the defensive industry. 

\vspace{-5pt}
\subsection{ESG investment and global political situation}\label{ESG investment and global political situation}
\vspace{-10pt}
\hspace{10pt} Contrary to the narrative in the article, despite an explosive increase in investment in ESG-related products, we are still witnessing a pullback from equity investment firms and their respective ESG factor division in the past two years. As well as the attack from political parties that deemed ESG-related investment anti-oil and are specifically targeting the fossil fuel industry. Looking at recent changes regarding ESG investment, early this year, BlackRock off 600 people, mainly from their ESG division \cite{Fox}. The same thing happened last year, when Morningstar reduced 10\% to 12\% of the global headcount of its Sustainalytics unit that provides research and ratings on ESG factors\cite{R}. Vanguard is the second-largest investment asset management in the world, withdrawn from the Net Zero Asset Managers initiative which was the the world’s largest climate-finance alliance\cite{Vanguard}. The ambiguous attitude of global investment/asset manager firms towards ESG and climate-related initiatives, and the potential political struggle bestowed upon the corporations may create an unwanted economic struggle for ESG factor-related investment and ESG-disclosed companies. 

\vspace{-5pt}
\subsection{Mixed empirical evidence of ESG factor's impact on corporation value}\label{Mixed empirical evidence of ESG factor's impact on corporation value}
\vspace{-10pt}
\hspace{10pt} Undergoing projects that increase ESG-related factors will impose more costs. Undertaking projects to increase different ESG factors will stretch out the corporation's budget with no real focus. Here the three main hypotheses are that companies are compensated with lower cost of capital, lower cost of debt, and lower cost of equity to increase the range of future projects the corporation could invest in and thus increase the competitive advantage that potentially has higher economic growth than the industry competition. 

However, empirical evidence shows mixed results regarding these three hypotheses. In one of the studies (Yusong Chen and Jinyi Yue 2022)\cite{Yusong}, researchers look into the relationship between ESG performances and ESG disclosure of those signed up to the Climate Action 100+ initiative, where they find no
significant correlation between the ESG performances with Cost of Capital, Cost of Debt, or Cost of Equity. Upon decomposing the main ESG factors into respective main Environmental, Social, and Corporate Governance pillars, only the Social pillars exhibit a negative relationship across the cost of debt, capital, and equity. ESG disclosure consistently indicates that firms with higher ESG disclosure have a higher cost of capital and typically a higher cost of debt.

Another study in the Banking industry (Andrieș et al 2023)\cite{Andries} suggests that Banks benefit from incorporating ESG practices into financial decisions, and they benefit from having a lower cost of raising interest-bearing liabilities, as well as a reduced cost of attracting deposits. Also, the study indicates only large banks and those located in developed countries benefit from increased transparency in the reporting of ESG information in
terms of reduced funding costs.
\vspace{-5pt}


%-----------------------------------------Conclusion-----------------------------------
% \textcolor{red}{(I don't think we varied the temperature enough - Lucas)}
\section{Extension} \label{Extension}
\vspace{-10pt}
\hspace{10pt} There are three aspects the authors could extend upon the original paper. Firstly, the growth of the industry when the industry starts to embrace the Stakeholder model and invest more into ESG-related factors within the corporation, such as the Oil and Gas industry. Secondly, the scale of different companies within the same industry and preferably the same country. Expand upon how corporations with different scales take the approach to the commitment of the Stakeholder model and how the corporation performs before and after the transition. Lastly, the authors should shed more insight into ESG score rating agencies and their relative relationship between their customers and market investors. Make a direct comparison of ESG rating agencies and credit rating agencies. What are the supposed differences between Mortgage Back Security rating and the corporation ESG score rating, and how differently do the regulators behave now, versus how they behaved pre-2008 financial crisis? 

\newpage
\printbibliography %Prints bibliography

\end{document}
